% --- Archon physics formalization source begin ---
% archon:physics
% archon:covers PhyXMiniProblems/problem_phyx_mini_0399.lean
% archon:source-report reports/phyx_mini/problem_phyx_mini_0399.source.json

\chapter{Physics problem phyx\_mini\_0399}
\label{ch:PhyXMiniProblems_problem_phyx_mini_0399}

\paragraph{Problem source.}
\#\# Physical scenario

Air in an automobile tire is initially at \$-10\textasciicircum{}\{\textbackslash{}circ\}C\$ and 190 kPa. After the automobile is driven awhile, the temperature rises to \$10\textasciicircum{}\{\textbackslash{}circ\}C\$.

\#\# Question

Find the new pressure. You must make one assumption on your own.

\#\# Answer choices

- A: A: 4MPa

- B: B: 204.4kPa

- C: C: 154kPa

- D: D: 10.2kPa

\#\# Figure caption (auxiliary; use the image as primary evidence)

The image shows a detailed illustration of a car tire. The tire has a visible tread pattern on its surface and is mounted on a wheel rim with a distinctive design. Attached to the wheel is a valve stem, which is indicated by an arrow labeled "Air P." This label suggests that air pressure is related to the valve. The rim is depicted with a blue hue, adding contrast to the black tire and emphasizing the structural components.

\#\# Dataset metadata

Ideal Gases and Kinetic Theory; Physical Model Grounding Reasoning, Implicit Condition Reasoning

\paragraph{Recorded answer/context.}
B: B: 204.4kPa

\paragraph{Figure/image path.}
/root/proposal\_for\_physic/science-mango/phyx\_mini\_run/phyx\_data/test\_image/399.png

\paragraph{Formalization target.}
create a compiling Lean file with sorry bodies at `PhyXMiniProblems/problem\_phyx\_mini\_0399.lean`. The Lean declarations must preserve the physical quantities, dimensions or dimensional roles, figure labels, governing-law hypotheses, and final relation expressed by this problem.
Use Mathlib/Physlib names found through LeanExplore where available. If a physics API is missing, introduce faithful local abstractions rather than scalar placeholder aliases.

\begin{theorem}[Physics formalization target]
\label{thm:physics:phyx_mini_0399:target}
\lean{PhyXMiniProblems.ProblemPhyXMini0399.new_tire_pressure_is_answer_B}
\uses{def:physics:phyx-mini-0399:phyxminiproblems-problemphyxmini0399-pressureinkilopascals, def:physics:phyx-mini-0399:phyxminiproblems-problemphyxmini0399-tireheatingsetup, def:physics:phyx-mini-0399:phyxminiproblems-problemphyxmini0399-matchesproblemstatement, def:physics:phyx-mini-0399:phyxminiproblems-problemphyxmini0399-matchesprimarytirefigure, def:physics:phyx-mini-0399:phyxminiproblems-problemphyxmini0399-hasphysicaltireairparameters, def:physics:phyx-mini-0399:phyxminiproblems-problemphyxmini0399-assumesconstanttirevolume, def:physics:phyx-mini-0399:phyxminiproblems-problemphyxmini0399-satisfiesfixedsampleidealgaslaw, def:physics:phyx-mini-0399:phyxminiproblems-problemphyxmini0399-recordeddatasetanswer, def:physics:phyx-mini-0399:phyxminiproblems-problemphyxmini0399-roundstodisplayedtenthkilopascal, def:physics:phyx-mini-0399:phyxminiproblems-problemphyxmini0399-isuniqueclosestdisplayedpressure}
The assigned autoformalize agent should translate this physics problem into Lean declarations in the covered file, with theorem and lemma proof bodies written as `by sorry`.
\end{theorem}
\begin{proof}
This is an autoformalization task, not a proof task. Produce statements that can later be proved by the physics prover without weakening the physical model.
\end{proof}

% --- Archon declaration topology begin ---
\section{Declaration topology}
\begin{definition}[Volume Quantity]
  \label{def:physics:phyx-mini-0399:phyxminiproblems-problemphyxmini0399-volumequantity}
  \lean{PhyXMiniProblems.ProblemPhyXMini0399.VolumeQuantity}
  Physical volume of the tire interior, with dimension L³
\end{definition}
\begin{proof}
  This is the declared model definition of Volume Quantity.
\end{proof}
\begin{definition}[pressure In Pascals]
  \label{def:physics:phyx-mini-0399:phyxminiproblems-problemphyxmini0399-pressureinpascals}
  \lean{PhyXMiniProblems.ProblemPhyXMini0399.pressureInPascals}
  Read a physical pressure in coherent SI pascals
\end{definition}
\begin{proof}
  This is the declared model definition of pressure In Pascals.
\end{proof}
\begin{definition}[pressure In Kilopascals]
  \label{def:physics:phyx-mini-0399:phyxminiproblems-problemphyxmini0399-pressureinkilopascals}
  \lean{PhyXMiniProblems.ProblemPhyXMini0399.pressureInKilopascals}
  \uses{def:physics:phyx-mini-0399:phyxminiproblems-problemphyxmini0399-pressureinpascals}
  Read a physical pressure in kilopascals
\end{definition}
\begin{proof}
  This is the declared model definition of pressure In Kilopascals.
\end{proof}
\begin{definition}[volume In Cubic Meters]
  \label{def:physics:phyx-mini-0399:phyxminiproblems-problemphyxmini0399-volumeincubicmeters}
  \lean{PhyXMiniProblems.ProblemPhyXMini0399.volumeInCubicMeters}
  \uses{def:physics:phyx-mini-0399:phyxminiproblems-problemphyxmini0399-volumequantity}
  Read a physical volume in coherent SI cubic metres
\end{definition}
\begin{proof}
  This is the declared model definition of volume In Cubic Meters.
\end{proof}
\begin{definition}[temperature In Kelvins]
  \label{def:physics:phyx-mini-0399:phyxminiproblems-problemphyxmini0399-temperatureinkelvins}
  \lean{PhyXMiniProblems.ProblemPhyXMini0399.temperatureInKelvins}
  Convert a stored Physlib absolute temperature to a kelvin readout
\end{definition}
\begin{proof}
  This is the declared model definition of temperature In Kelvins.
\end{proof}
\begin{definition}[celsius Zero In Kelvins]
  \label{def:physics:phyx-mini-0399:phyxminiproblems-problemphyxmini0399-celsiuszeroinkelvins}
  \lean{PhyXMiniProblems.ProblemPhyXMini0399.celsiusZeroInKelvins}
  The exact Celsius-zero offset 273.15 K
\end{definition}
\begin{proof}
  This is the declared model definition of celsius Zero In Kelvins.
\end{proof}
\begin{definition}[temperature In Degrees Celsius]
  \label{def:physics:phyx-mini-0399:phyxminiproblems-problemphyxmini0399-temperatureindegreescelsius}
  \lean{PhyXMiniProblems.ProblemPhyXMini0399.temperatureInDegreesCelsius}
  \uses{def:physics:phyx-mini-0399:phyxminiproblems-problemphyxmini0399-temperatureinkelvins, def:physics:phyx-mini-0399:phyxminiproblems-problemphyxmini0399-celsiuszeroinkelvins}
  Celsius readout obtained from the physical absolute temperature
\end{definition}
\begin{proof}
  This is the declared model definition of temperature In Degrees Celsius.
\end{proof}
\begin{definition}[Driving Stage]
  \label{def:physics:phyx-mini-0399:phyxminiproblems-problemphyxmini0399-drivingstage}
  \lean{PhyXMiniProblems.ProblemPhyXMini0399.DrivingStage}
  The tire-air states before and after the automobile is driven
\end{definition}
\begin{proof}
  This is the declared model definition of Driving Stage.
\end{proof}
\begin{definition}[Gas Model]
  \label{def:physics:phyx-mini-0399:phyxminiproblems-problemphyxmini0399-gasmodel}
  \lean{PhyXMiniProblems.ProblemPhyXMini0399.GasModel}
  Equation-of-state model assigned to the tire air
\end{definition}
\begin{proof}
  This is the declared model definition of Gas Model.
\end{proof}
\begin{definition}[Tire Air State]
  \label{def:physics:phyx-mini-0399:phyxminiproblems-problemphyxmini0399-tireairstate}
  \lean{PhyXMiniProblems.ProblemPhyXMini0399.TireAirState}
  \uses{def:physics:phyx-mini-0399:phyxminiproblems-problemphyxmini0399-volumequantity}
  A thermodynamic state of the air inside the tire
\end{definition}
\begin{proof}
  This is the declared model definition of Tire Air State.
\end{proof}
\begin{definition}[Tire Air Sample]
  \label{def:physics:phyx-mini-0399:phyxminiproblems-problemphyxmini0399-tireairsample}
  \lean{PhyXMiniProblems.ProblemPhyXMini0399.TireAirSample}
  \uses{def:physics:phyx-mini-0399:phyxminiproblems-problemphyxmini0399-gasmodel}
  The same physical air sample occurs at both stages. The scalar field is an explicit SI readout of the sample-dependent quantity nR, not a replacement type for the gas sample
\end{definition}
\begin{proof}
  This is the declared model definition of Tire Air Sample.
\end{proof}
\begin{definition}[Tire Figure Component]
  \label{def:physics:phyx-mini-0399:phyxminiproblems-problemphyxmini0399-tirefigurecomponent}
  \lean{PhyXMiniProblems.ProblemPhyXMini0399.TireFigureComponent}
  Components visibly present in the supplied tire illustration
\end{definition}
\begin{proof}
  This is the declared model definition of Tire Figure Component.
\end{proof}
\begin{definition}[Figure Observable]
  \label{def:physics:phyx-mini-0399:phyxminiproblems-problemphyxmini0399-figureobservable}
  \lean{PhyXMiniProblems.ProblemPhyXMini0399.FigureObservable}
  Physical observable denoted by a label in the supplied illustration
\end{definition}
\begin{proof}
  This is the declared model definition of Figure Observable.
\end{proof}
\begin{definition}[Tire Pressure Figure]
  \label{def:physics:phyx-mini-0399:phyxminiproblems-problemphyxmini0399-tirepressurefigure}
  \lean{PhyXMiniProblems.ProblemPhyXMini0399.TirePressureFigure}
  \uses{def:physics:phyx-mini-0399:phyxminiproblems-problemphyxmini0399-tirefigurecomponent, def:physics:phyx-mini-0399:phyxminiproblems-problemphyxmini0399-figureobservable}
  Qualitative data transcribed from the bitmap. The label and arrow identify the pressure observable at the valve stem; the bitmap supplies no pressure number
\end{definition}
\begin{proof}
  This is the declared model definition of Tire Pressure Figure.
\end{proof}
\begin{definition}[Tire Heating Setup]
  \label{def:physics:phyx-mini-0399:phyxminiproblems-problemphyxmini0399-tireheatingsetup}
  \lean{PhyXMiniProblems.ProblemPhyXMini0399.TireHeatingSetup}
  \uses{def:physics:phyx-mini-0399:phyxminiproblems-problemphyxmini0399-drivingstage, def:physics:phyx-mini-0399:phyxminiproblems-problemphyxmini0399-tireairstate, def:physics:phyx-mini-0399:phyxminiproblems-problemphyxmini0399-tireairsample, def:physics:phyx-mini-0399:phyxminiproblems-problemphyxmini0399-tirepressurefigure}
  The fixed gas sample, its two states, and the supplied illustration
\end{definition}
\begin{proof}
  This is the declared model definition of Tire Heating Setup.
\end{proof}
\begin{definition}[Matches Problem Statement]
  \label{def:physics:phyx-mini-0399:phyxminiproblems-problemphyxmini0399-matchesproblemstatement}
  \lean{PhyXMiniProblems.ProblemPhyXMini0399.MatchesProblemStatement}
  \uses{def:physics:phyx-mini-0399:phyxminiproblems-problemphyxmini0399-pressureinkilopascals, def:physics:phyx-mini-0399:phyxminiproblems-problemphyxmini0399-temperatureindegreescelsius, def:physics:phyx-mini-0399:phyxminiproblems-problemphyxmini0399-tireheatingsetup}
  Numerical data stated in the prose. The requested final pressure deliberately does not occur in this structure
\end{definition}
\begin{proof}
  This is the declared model definition of Matches Problem Statement.
\end{proof}
\begin{definition}[Matches Primary Tire Figure]
  \label{def:physics:phyx-mini-0399:phyxminiproblems-problemphyxmini0399-matchesprimarytirefigure}
  \lean{PhyXMiniProblems.ProblemPhyXMini0399.MatchesPrimaryTireFigure}
  \uses{def:physics:phyx-mini-0399:phyxminiproblems-problemphyxmini0399-tireheatingsetup}
  Primary-image evidence: the tire, blue wheel rim, and valve stem are visible, and the arrow labelled Air P points to the valve stem to denote air pressure. There is no numerical figure readout to assume
\end{definition}
\begin{proof}
  This is the declared model definition of Matches Primary Tire Figure.
\end{proof}
\begin{definition}[Has Physical Tire Air Parameters]
  \label{def:physics:phyx-mini-0399:phyxminiproblems-problemphyxmini0399-hasphysicaltireairparameters}
  \lean{PhyXMiniProblems.ProblemPhyXMini0399.HasPhysicalTireAirParameters}
  \uses{def:physics:phyx-mini-0399:phyxminiproblems-problemphyxmini0399-pressureinpascals, def:physics:phyx-mini-0399:phyxminiproblems-problemphyxmini0399-volumeincubicmeters, def:physics:phyx-mini-0399:phyxminiproblems-problemphyxmini0399-temperatureinkelvins, def:physics:phyx-mini-0399:phyxminiproblems-problemphyxmini0399-drivingstage, def:physics:phyx-mini-0399:phyxminiproblems-problemphyxmini0399-tireheatingsetup}
  Positivity conditions selecting physically meaningful tire-air states
\end{definition}
\begin{proof}
  This is the declared model definition of Has Physical Tire Air Parameters.
\end{proof}
\begin{definition}[Assumes Constant Tire Volume]
  \label{def:physics:phyx-mini-0399:phyxminiproblems-problemphyxmini0399-assumesconstanttirevolume}
  \lean{PhyXMiniProblems.ProblemPhyXMini0399.AssumesConstantTireVolume}
  \uses{def:physics:phyx-mini-0399:phyxminiproblems-problemphyxmini0399-tireheatingsetup}
  The one modeling assumption requested by the problem: tire deformation is neglected, so the air occupies the same volume before and after driving
\end{definition}
\begin{proof}
  This is the declared model definition of Assumes Constant Tire Volume.
\end{proof}
\begin{definition}[Satisfies Fixed Sample Ideal Gas Law]
  \label{def:physics:phyx-mini-0399:phyxminiproblems-problemphyxmini0399-satisfiesfixedsampleidealgaslaw}
  \lean{PhyXMiniProblems.ProblemPhyXMini0399.SatisfiesFixedSampleIdealGasLaw}
  \uses{def:physics:phyx-mini-0399:phyxminiproblems-problemphyxmini0399-pressureinpascals, def:physics:phyx-mini-0399:phyxminiproblems-problemphyxmini0399-volumeincubicmeters, def:physics:phyx-mini-0399:phyxminiproblems-problemphyxmini0399-temperatureinkelvins, def:physics:phyx-mini-0399:phyxminiproblems-problemphyxmini0399-drivingstage, def:physics:phyx-mini-0399:phyxminiproblems-problemphyxmini0399-tireheatingsetup}
  The macroscopic ideal-gas equation pV = (nR)T, in coherent SI readouts, at both states of the same air sample. This governing law contains no numerical final pressure
\end{definition}
\begin{proof}
  This is the declared model definition of Satisfies Fixed Sample Ideal Gas Law.
\end{proof}
\begin{definition}[Answer Choice]
  \label{def:physics:phyx-mini-0399:phyxminiproblems-problemphyxmini0399-answerchoice}
  \lean{PhyXMiniProblems.ProblemPhyXMini0399.AnswerChoice}
  Labels of the four multiple-choice answers
\end{definition}
\begin{proof}
  This is the declared model definition of Answer Choice.
\end{proof}
\begin{definition}[displayed Pressure Kilopascals]
  \label{def:physics:phyx-mini-0399:phyxminiproblems-problemphyxmini0399-displayedpressurekilopascals}
  \lean{PhyXMiniProblems.ProblemPhyXMini0399.displayedPressureKilopascals}
  \uses{def:physics:phyx-mini-0399:phyxminiproblems-problemphyxmini0399-answerchoice}
  Pressure displayed beside each answer, normalized to kilopascals. In particular, choice A's 4 MPa is represented as 4000 kPa
\end{definition}
\begin{proof}
  This is the declared model definition of displayed Pressure Kilopascals.
\end{proof}
\begin{definition}[recorded Dataset Answer]
  \label{def:physics:phyx-mini-0399:phyxminiproblems-problemphyxmini0399-recordeddatasetanswer}
  \lean{PhyXMiniProblems.ProblemPhyXMini0399.recordedDatasetAnswer}
  \uses{def:physics:phyx-mini-0399:phyxminiproblems-problemphyxmini0399-answerchoice}
  Answer label recorded in the supplied dataset metadata
\end{definition}
\begin{proof}
  This is the declared model definition of recorded Dataset Answer.
\end{proof}
\begin{definition}[Rounds To Displayed Tenth Kilopascal]
  \label{def:physics:phyx-mini-0399:phyxminiproblems-problemphyxmini0399-roundstodisplayedtenthkilopascal}
  \lean{PhyXMiniProblems.ProblemPhyXMini0399.RoundsToDisplayedTenthKilopascal}
  \uses{def:physics:phyx-mini-0399:phyxminiproblems-problemphyxmini0399-answerchoice, def:physics:phyx-mini-0399:phyxminiproblems-problemphyxmini0399-displayedpressurekilopascals}
  The physical pressure rounds to the displayed tenth of a kilopascal
\end{definition}
\begin{proof}
  This is the declared model definition of Rounds To Displayed Tenth Kilopascal.
\end{proof}
\begin{definition}[Is Unique Closest Displayed Pressure]
  \label{def:physics:phyx-mini-0399:phyxminiproblems-problemphyxmini0399-isuniqueclosestdisplayedpressure}
  \lean{PhyXMiniProblems.ProblemPhyXMini0399.IsUniqueClosestDisplayedPressure}
  \uses{def:physics:phyx-mini-0399:phyxminiproblems-problemphyxmini0399-answerchoice, def:physics:phyx-mini-0399:phyxminiproblems-problemphyxmini0399-displayedpressurekilopascals}
  The selected answer is closer than every alternative pressure display
\end{definition}
\begin{proof}
  This is the declared model definition of Is Unique Closest Displayed Pressure.
\end{proof}
% --- Archon declaration topology end ---
% --- Archon physics formalization source end ---
